\oldpage[I]{81}
\begin{env}[1.1.3]
\label{I.1.1.3-fr}
Pour toute partie $Y$ de $X$, nous désignerons par $\mathfrak{j}(Y)$
l'ensemble des $f\in A$ tels que $f(y)=0$ pour tout $y\in Y$~; il revient au
même de dire que $\mathfrak{j}(Y)$ est l'intersection des idéaux premiers
$\mathfrak{j}_y$ pour $y\in Y$. Il est clair que la relation $Y\subset Y'$
entraîne $\mathfrak{j}(Y)\supset\mathfrak{j}(Y')$ et que l'on a
\begin{equation}
  \mathfrak{j}\left(\bigcup_\lambda Y_\lambda\right)
    =\bigcap_\lambda\mathfrak{j}(Y_\lambda)
  \tag{1.1.3.1}\label{I.1.1.3.1-fr}
\end{equation}
pour toute famille $(Y_\lambda)$ de parties de $X$. Enfin on a
\begin{equation}
  \mathfrak{j}(\{x\})=\mathfrak{j}_x.
  \tag{1.1.3.2}\label{I.1.1.3.2-fr}
\end{equation}
\end{env}

\begin{proposition}[1.1.4]
\label{I.1.1.4-fr}
\begin{enumerate}
  \item[\rm(i)] Pour toute partie $E$ de $A$, on a
  $\mathfrak{j}(V(E))=\mathfrak{r}(E)$.
  \item[\rm(ii)] Pour toute partie $Y$ de $X$,
  $V(\mathfrak{j}(Y))=\overline{Y}$, adhérence de $Y$ dans $X$.
\end{enumerate}
\end{proposition}

(i) est conséquence immédiate des définitions et de
(\hyperref[I.1.1.1.1-fr]{1.1.1.1})~; d'autre part,
$V(\mathfrak{j}(Y))$ est fermé et contient $Y$~; inversement, si
$Y\subset V(E)$, on a $f(y)=0$ pour tout $f\in E$ et tout $y\in Y$, donc
$E\subset\mathfrak{j}(Y)$, $V(E)\supset V(\mathfrak{j}(Y))$, ce qui prouve
(ii).

\begin{corollary}[1.1.5]
\label{I.1.1.5-fr}
Les parties fermées de $X=\operatorname{Spec}(A)$ et les idéaux de $A$ égaux
à leurs racines (autrement dit les intersections d'idéaux premiers) se
correspondent biunivoquement par les applications décroissantes
$Y\mapsto\mathfrak{j}(Y)$, $\mathfrak{a}\mapsto V(\mathfrak{a})$~; à la
réunion $Y_1\cup Y_2$ de deux parties fermées correspond
$\mathfrak{j}(Y_1)\cap\mathfrak{j}(Y_2)$, et à l'intersection d'une famille
quelconque $(Y_\lambda)$ de parties fermées correspond la racine de la somme
des $\mathfrak{j}(Y_\lambda)$.
\end{corollary}

\begin{corollary}[1.1.6]
\label{I.1.1.6-fr}
Si $A$ est un anneau noethérien, $X=\operatorname{Spec}(A)$ est un espace
noethérien.
\end{corollary}

On notera que la réciproque de ce corollaire est inexacte, comme le montre
l'exemple d'un anneau intègre non noethérien ayant un seul idéal premier
$\neq\{0\}$, par exemple un anneau de valuation non discrète de rang 1.

Comme exemple d'anneau $A$ dont le spectre n'est pas un espace noethérien, on
peut signaler l'anneau $\mathcal{C}(Y)$ des fonctions continues numériques sur
un espace compact infini $Y$~; on sait qu'en tant qu'ensemble, $Y$ s'identifie
à l'ensemble des idéaux maximaux de $A$, et il est facile de voir que la
topologie induite sur $Y$ par celle de $X=\operatorname{Spec}(A)$ est la
topologie initiale de $Y$. Comme $Y$ n'est pas un espace noethérien, il en est
donc de même de $X$.

\begin{corollary}[1.1.7]
\label{I.1.1.7-fr}
Pour tout $x\in X$, l'adhérence de $\{x\}$ est l'ensemble des $y\in X$ tels
que $\mathfrak{j}_x\subset\mathfrak{j}_y$. Pour que $\{x\}$ soit fermé, il
faut et il suffit que $\mathfrak{j}_x$ soit maximal.
\end{corollary}

\begin{corollary}[1.1.8]
\label{I.1.1.8-fr}
L'espace $X=\operatorname{Spec}(A)$ est un espace de Kolmogoroff.
\end{corollary}

En effet, si $x$, $y$ sont deux points distincts de $X$, on a, soit
$\mathfrak{j}_x\not\subset\mathfrak{j}_y$, soit
$\mathfrak{j}_y\not\subset\mathfrak{j}_x$, donc l'un des points $x$, $y$
n'appartient pas à l'adhérence de l'autre.

\begin{env}[1.1.9]
\label{I.1.1.9-fr}
D'après la prop. (\hyperref[I.1.1.2-fr]{1.1.2}, (iv)), pour deux éléments
$f$, $g$ de $A$, on a
\begin{equation}
  D(fg)=D(f)\cap D(g).
  \tag{1.1.9.1}\label{I.1.1.9.1-fr}
\end{equation}
Notons aussi que la relation $D(f)=D(g)$ signifie, d'après la prop.
(\hyperref[I.1.1.4-fr]{1.1.4}, (i)) et la prop.
(\hyperref[I.1.1.2-fr]{1.1.2}, (v)) que
$\mathfrak{r}(f)=\mathfrak{r}(g)$, ou encore que les idéaux premiers minimaux
contenant $(f)$ et $(g)$ sont les mêmes~; en particulier, il en est ainsi
lorsque $f=ug$, où $u$ est inversible.
\end{env}

